% ============================================================
%  第二卷 第一章 相对性原理（§1–§7；习题仅在 §6、§7）
%  底本：高教社中文版第8版；OCR 错误已修正，脚注文本待从纸版补录
%  已修正的 OCR 错误（行号=full_merged.md）：
%  - 行373 段末乱码"Xa"，删
%  - 行391 "足够精确"后补句点（原书 p.2 有）
%  - 行403 "才有意义"后补句点；行538 "类时的"后补句点；
%    行607 "固有时"后补句点；行617 "∫ds"后补句点
%  - 行414 半角逗号"，"→全角；行683 半角分号";"→全角
%  - 行458 原书衍字："以 x,y,z 和积 ct 标记"→"和 ct 标记"（原书 p.5 排印衍字）
%  - 行460 \mathbf{d}s′→ds′（非粗体）
%  - 行524 "l′²₁₂=0"后补句点；行560 "$c_\cdot$"→"$c$."
%  - 行575、747 "引人"→"引入"
%  - 行621 句末补句点；行631、635、771 段末补句点
%  - 行643 "=x′²"后补句点；行679 "（4.3)得到"前补句点
%  - 行759 OCR 乱码"$\sf F$"→"下"（极限情形下）
%  - 行897、905 混合分量第二处 "$A_k$"→"$A_i{}^k$"（原书 p.17）
%  - 行917 "$A^iB_i$"后补句点；行915 "缩并"后补句点
%  - 行928 (6.4) OCR 乱码"av=1"→"当"（分段函数条件）
%  - 行936 (6.5) 矩阵后补句点；行1001 段末补句点
%  - 行963 乱码"\bar{j}^\pi 格地说"→"严格地说"
%  - 行970 四公式行列式第3行指标 OCR 错乱，按原书改 δ_p^l,δ_r^l,δ_s^l,δ_t^l
%  - 行980 乱码"\Re\underline{\mathbb{E}}"→"其中"
%  - 行992、998 p、a 补粗体；行1023 dr→\boldsymbol{r}
%  - 行1045 "等于并“垂直”于面元 df^{ik}"后补逗号
%  - 行1120 (6.18) OCR 丢箭头，补"→"
%  - 行1195 "求 u^2,u^3,u^4"→"u^0,u^2,u^3"（原书 p.25 排印错误，分量指标只有0–3）
%  - 行1238 乱码"{\mathfrak X}"→"或"；行1241、1247 粗体 t=0,v=0,x=0 → 正体
%  - 行1253 "wt∞"→"wt→∞"、"恒定值C"→"恒定值 c"（原书 p.26）
%  - 行1261 按原书 p.26 补"ln"：(c/w)ln(2wt/c)
%  排版说明：本章 §3 首段行内根式不可断行，为消除 overfull，本章内
%  \emergencystretch=3em（仅影响原本会溢出的行，不改变文字）
%  遗留问题：
%  - 脚注已转录（2026-10-10 脚注代理自原书扫描页补录，全章共11处）。
%    其中原书 p.19（p-033.png）脚注①"我们给出下列公式以备参考…"所含 4 组
%    公式此前被误置于正文（原 563–579 行两个 equation*），现已按原书移入脚注
%  - (6.4) 分段条件"当 i=k／当 i≠k"之"当"为按文意补（OCR 乱码）
% ============================================================
\chapter{相对性原理}
\emergencystretch 3em

\section{相互作用的传播速度}

为了描述自然界中所发生的过程，必须有一个所谓参考系.参考系应理解为一个坐标系和固定在这个坐标系里的钟.坐标系用来刻画一个粒子在空间的位置，钟用来指示时间.

有这样一类参考系，在其中，一个自由运动物体，即一个无外力作用于其上的运动物体，是以恒定速度行进的.这类参考系叫做惯性系.

如果两个参考系彼此相对作匀速直线运动，而其中的一个又是惯性系，那么，另外一个显然也是惯性系（在这个参考系中每一个自由运动也将是匀速直线运动）.因此，我们可以有任意多个惯性参考系，它们彼此相对作匀速直线运动.

实验表明，所谓相对性原理是有效的.按照这个原理，所有的自然定律在所有惯性参考系中都是相同的.换句话说，表示自然定律的方程对于由一个惯性系到另一个惯性系的时间与坐标的各种变换来说是不变的.这就是说，描述自然界定律的方程，如用不同的惯性参考系的坐标与时间写出来，将有同样的形式.

粒子间的相互作用在普通力学中由相互作用势能来描述，相互作用势能是相互作用的粒子的坐标的函数.很容易看出，这种描述相互作用的方式，包含着一个假定，即假定相互作用是瞬时传播的.事实上，按照上面的说法，每一个粒子在某一瞬时受到其他各粒子的作用力，仅与那些粒子在该瞬时的位置有关.在这些相互作用的粒子中，如果有一个粒子改变了位置，立刻就会影响到其他各粒子.

然而，实验表明，瞬时的相互作用在自然界中是不存在的.因此，基于相互作用的瞬时传播概念的力学本身就含有某些不准确性.实际上，如果相互作用的物体中的一个发生任何变动，仅仅在过了某段时间以后才能影响到其他物体.只有在这段时间以后，由最初变动所产生的物理过程才开始在第二个物体上发生.用这段时间除两个物体间的距离，就得到相互作用的传播速度.

我们要注意，这个速度，严格地说，应该称为相互作用的最大传播速度.这个速度仅仅决定某一物体的变动开始在第二个物体上表现出来所需要的时间间隔.显然，相互作用的最大传播速度的存在，同时也就暗示着，在自然界中，物体运动的速度一般不可能大于这个速度.事实上，假若真的有这种运动存在，那么我们就可以利用这种运动实现一个相互作用，其传播速度比上面所说的最大传播速度还要大.

从一个粒子向另一个粒子传播的相互作用往往叫做“信号”，它由第一个粒子发出，将第一个粒子所经历的变化“通知”第二个粒子.因此相互作用的传播速度称为信号速度.

值得注意的是，由相对性原理可以推断相互作用的传播速度在所有惯性参考系中都是一样的.因此，相互作用的传播速度是一个普适常数.

以后我们将要证明，这个恒定速度就是光在真空中的速度.我们通常用字母 $c$ 来代表光速，其值等于
\begin{equation}
  c=2.998\times 10^{10}\ \mathrm{cm/s}.
\end{equation}

这个速度很高，这可以解释经典力学为何在大多数情况下都足够精确.我们有机会遇到的各种速度通常都比光速小得多，以至假设光速为无限大，对结果的精确性并无实质上的影响.

把相对性原理同相互作用传播速度的有限性结合起来，就是爱因斯坦的相对性原理（爱因斯坦在1905年提出这个原理），它不同于伽利略的相对性原理，伽利略的相对性原理基于无限大的相互作用传播速度.

以爱因斯坦的相对性原理（以后我们通常简称它为相对性原理）为基础的力学，称为相对论力学.在运动物体的速度远小于光速的极限情形下，我们就可以略去传播速度的有限性对于运动的影响.这样一来，相对论力学就变为通常的力学了，通常的力学基于相互作用瞬时传播这一假定；这种力学称为牛顿力学或经典力学.在相对论力学的公式中，取 $c\to\infty$ 极限，就可由相对论力学在形式上过渡到经典力学.

在经典力学中，距离已经是相对的，就是说，不同事件的空间关系依赖于描述这些事件所用的参考系.所以，说两件不同时的事件发生在空间同一点上，或者更普遍而言，说两件不同时的事件发生在彼此间有一定距离的两点上，只有当我们指明了我们所用的是哪一个参考系时才有意义.

另一方面，在经典力学中，时间是绝对的；换句话说，经典力学假定时间的特性与参考系无关，对所有参考系来说，时间只有一个.这就是说，假如对于某一个观察者来说，有两个现象是同时发生的，那么，对所有其他观察者来说，这两个现象也是同时发生的.更普遍而言，两个给定事件发生的时间间隔在一切参考系中必须一样.

然而，很容易证明，绝对时间的概念是与爱因斯坦的相对性原理完全冲突的.为了说明这一点，我们只需回忆一下，在以绝对时间的概念为基础的经典力学中，速度合成的通用法则是有效的.按照这个法则，复合运动的合速度简单地等于组成这个运动的各个速度的（矢量）和.这个法则既然是普遍适用的，就应该可以应用于相互作用的传播.由此可以推出，传播速度在不同的惯性参考系中必定是不同的，这就与相对性原理冲突了.但是，实验完全证实了相对性原理.在1881年迈克耳孙首次测量的结果显示，光速与其传播方向并无关系；然而，按照经典力学，光速在与地球运动方向相同的方向上，应该比在与地球运动方向相反的方向上为小.

因此，相对性原理导出一个结果，即时间不是绝对的.在不同的参考系中，时间的流逝也是不同的.所以，“两个不同的事件之间有一定的时间间隔”这样的陈述，仅在肯定地指明了所应用的是哪一个参考系的情况下才有意义.特别是，在某一个参考系内同时发生的事件，对另一个参考系来说并不是同时的.

为了弄清楚这个概念，我们先考虑下面的简单例子.

我们来研究两个惯性参考系 $K$ 和 $K'$，其坐标轴分别为 $xyz$ 及 $x'y'z'$，而 $K'$ 则相对于 $K$ 沿 $x$ 和 $x'$ 轴向右运动（图1）.

\par\nopagebreak\noindent\begin{minipage}{\linewidth}\centering\includegraphics[width=5.5cm]{fig1.pdf}\\[0.3em]{\small 图 1}\end{minipage}\par\nopagebreak

设信号从 $x'$ 轴上某一点 $A$ 向两个相反的方向发出.既然信号在 $K'$ 系中的传播速度，正如在所有惯性系中一样，（在两个方向上）都等于 $c$，那么，它就会（在 $K'$ 系里）同时到达与 $A$ 等距离的两点 $B$ 及 $C$.

但是，很显然，同样的两事件（信号到达 $B$ 及 $C$ 两点），对于在 $K$ 系内的观察者来说，绝不是同时的.实际上，按照相对性原理，信号相对于 $K$ 系的速度也等于 $c$，并且因为 $B$ 点对于 $K$ 系而言，是对着向它发出的信号移动，而 $C$ 点则背离（由 $A$ 向 $C$ 发出的）信号移动，所以在 $K$ 系中，信号到达 $B$ 点要比到达 $C$ 点为早.

因此，爱因斯坦的相对性原理使基本物理概念发生了极深刻的和根本的改变.由我们日常生活经验所导出的空间和时间的概念仅仅是近似的，因为我们日常生活所遇到的速度，都比光速小得多.

\section{间隔}

以后我们常常要用事件这一概念.一个事件是由其发生的地点及其发生的时间来描述的.因此，在某一实物粒子上所发生的事件可由粒子的三个坐标及事件发生的时间来决定.

为表述便利起见，使用一个假想的四维空间往往是很实用的.在这个四维空间的四个轴中，三个用来刻画位置坐标，一个用来标示时间.在这个空间内，事件可用点来表示，这个点称为世界点.在这个假想的四维空间内，每一个粒子都对应于一条线，称为世界线.这条线上的各点决定了粒子在所有时刻的坐标.很容易证明，与一个作匀速直线运动的粒子相对应的世界线是一条直线.

现在我们用数学形式来表示光速不变原理.为此，我们考虑两个彼此以恒定速度作相对运动的参考系 $K$ 及 $K'$.这时我们选择 $x$ 轴与 $x'$ 轴重合，而 $y$ 和 $z$ 轴则分别与 $y'$ 和 $z'$ 轴平行，并以 $t$ 和 $t'$ 分别表示在 $K$ 和 $K'$ 参考系内的时间.

设第一个事件是：在 $K$ 系内的 $t_1$ 时刻从具有坐标 $x_1,y_1,z_1$（在同一参考系中）的点送出一个以光速传播的信号.我们就在 $K$ 系内观察这个信号的传播.再设第二个事件是：信号在 $t_2$ 时刻到达点 $x_2,y_2,z_2$，信号传播的速度既然是 $c$，所以它所经过的距离就是 $c(t_2-t_1)$.另一方面，这同一个距离又等于 $[(x_2-x_1)^2+(y_2-y_1)^2+(z_2-z_1)^2]^{1/2}$.因此，我们可以写出 $K$ 系内两个事件的坐标的关系：
\begin{equation}
  (x_2-x_1)^2+(y_2-y_1)^2+(z_2-z_1)^2-c^2(t_2-t_1)^2=0.
\end{equation}

同样两个事件，即该信号的传播，也可以在 $K'$ 系内观察：

设第一个事件在 $K'$ 内的坐标为 $x_1',y_1',z_1',t_1'$，而第二个事件则为 $x_2',y_2',z_2',t_2'$.按照光速不变原理，信号传播的速度在 $K'$ 系内与在 $K$ 系内相同，所以我们得到与（2.1）式相似的方程
\begin{equation}
  (x_2'-x_1')^2+(y_2'-y_1')^2+(z_2'-z_1')^2-c^2(t_2'-t_1')^2=0.
\end{equation}

假如 $x_1,y_1,z_1,t_1$ 及 $x_2,y_2,z_2,t_2$ 是任何两个事件的坐标，则
\begin{equation}
  s_{12}=[c^2(t_2-t_1)^2-(x_2-x_1)^2-(y_2-y_1)^2-(z_2-z_1)^2]^{1/2}
\end{equation}
称为这两个事件的间隔.

因此，由光速不变原理，我们可以断定，假如两个事件的间隔在某一个坐标系内为零，那么，它在所有其他坐标系内均为零.

如果两个事件彼此无限地接近，那么，其间隔 ds 将满足下面的方程：
\begin{equation}
  \mathrm{d}s^2=c^2\mathrm{d}t^2-\mathrm{d}x^2-\mathrm{d}y^2-\mathrm{d}z^2.
\end{equation}

从数学形式上看，表达式（2.3）和（2.4）容许我们把该间隔设想为四维空间内两点之间的距离（该空间的4个轴以 $x,y,z$ 和 $ct$ 标记）.但是，构成这个量的法则与普通几何的法则之间有一个根本区别：在构成间隔的平方时，沿不同轴的坐标差平方是以相异而非相同的运算符号求和的.\footnote{二次式（2.4）所描述的四维几何，是 H. 闵可夫斯基为相对论而引入的.这种几何称为伪欧几里得几何，以区别于普通的欧几里得几何.}

上面已经证明，如果在某一惯性系内 $\mathrm{d}s=0$，则在任一其他惯性系内也有 $\mathrm{d}s'=0$.此外，$\mathrm{d}s$ 与 $\mathrm{d}s'$ 为同阶的两个无穷小量.由以上两个情况可以得出结论，$\mathrm{d}s^2$ 与 $\mathrm{d}s'^2$ 彼此必须成比例：
\begin{equation*}
  \mathrm{d}s^2=a\,\mathrm{d}s'^2,
\end{equation*}
而且其中系数 $a$ 仅与两个惯性系的相对速度的绝对值有关.系数 $a$ 不可能与坐标或时间有关系，否则，空间的不同点及时间的不同时刻就不等价了，这是与时间及空间的均匀性相矛盾的.系数 $a$ 也不可能与惯性系的相对速度的方向有关，因为这就与空间的各向同性的性质相矛盾了.

考虑三个参考系 $K,K_1,K_2$，令 $V_1$ 和 $V_2$ 为 $K_1$ 和 $K_2$ 相对于 $K$ 的速度，则我们有：
\begin{equation*}
  \mathrm{d}s^2=a(V_1)\mathrm{d}s_1^2,\quad
  \mathrm{d}s^2=a(V_2)\mathrm{d}s_2^2.
\end{equation*}
类似地，我们可以写出
\begin{equation*}
  \mathrm{d}s_1^2=a(V_{12})\mathrm{d}s_2^2,
\end{equation*}
式中 $V_{12}$ 是 $K_2$ 相对于 $K_1$ 速度的绝对值.相互比较这些关系之后，我们发现必须有
\begin{equation}
  \frac{a(V_2)}{a(V_1)}=a(V_{12}).
\end{equation}

但是，$V_{12}$ 不仅依赖于矢量 $V_1$ 和 $V_2$ 的绝对值，而且依赖于它们之间的夹角.不过，这个角在（2.5）式左边并未出现.因此显而易见，这个公式只有当函数 $a(V)$ 为一常数时才成立，根据同一公式，该常数应等于1.

因此，
\begin{equation}
  \mathrm{d}s^2=\mathrm{d}s'^2,
\end{equation}
再从无限小间隔的相等得出有限间隔的相等：$s=s'$.

因此，我们得到一个很重要的结论：两个事件的间隔在所有惯性参考系里都是一样的，即当由一个惯性参考系变换到任何其他惯性参考系时，它是不变的.这个不变性就是光速不变的数学表示.

再次假设 $x_1,y_1,z_1,t_1$ 及 $x_2,y_2,z_2,t_2$ 是在某一个参考系 $K$ 内的两个事件的坐标.我们要问，是否有一个参考系 $K'$ 存在，在其中这两个事件在同一空间点发生？

我们采用下面的符号：
\begin{equation*}
  t_2-t_1=t_{12},\quad
  (x_2-x_1)^2+(y_2-y_1)^2+(z_2-z_1)^2=l_{12}^2.
\end{equation*}
于是，在 $K$ 系内，两个事件之间的间隔是：
\begin{equation*}
  s_{12}^2=c^2t_{12}^2-l_{12}^2,
\end{equation*}
而在 $K'$ 系内则是
\begin{equation*}
  s_{12}^{\prime 2}=c^2t_{12}^{\prime 2}-l_{12}^{\prime 2},
\end{equation*}
并且因为间隔的不变性，所以
\begin{equation*}
  c^2t_{12}^2-l_{12}^2=c^2t_{12}^{\prime 2}-l_{12}^{\prime 2}.
\end{equation*}

我们要求两个事件在 $K'$ 系中的同一点发生，即要求 $l_{12}^{\prime 2}=0$.这时
\begin{equation*}
  s_{12}^2=c^2t_{12}^2-l_{12}^2=c^2t_{12}^{\prime 2}>0.
\end{equation*}

因之，如果 $s_{12}^2>0$，即如果两个事件的间隔是实数的话，则具有我们所要求特性的参考系是存在的.实数间隔称为类时间隔.

因此，若两个事件的间隔是类时的，那么就有这样一个参考系存在，在其中两个事件发生于同一地点.在这个系统内，这两个事件的时间间隔等于
\begin{equation}
  t_{12}'=\frac{1}{c}\sqrt{c^2t_{12}^2-l_{12}^2}=\frac{s_{12}}{c}.
\end{equation}

若任何两个事件在同一物体上发生，那么，它们的间隔将永为类时的.事实上，因为物体运动的速度不可能大于 $c$，所以在两个事件之间，物体所行走的距离不可能大于 $ct_{12}$，因此我们永远有
\begin{equation*}
  l_{12}<ct_{12}.
\end{equation*}

现在我们再问，能否找到一个参考系，在其中这两个事件会同时发生？同上面一样，我们在 $K$ 及 $K'$ 两个参考系中有 $c^2t_{12}^2-l_{12}^2=c^2t_{12}^{\prime 2}-l_{12}^{\prime 2}$.我们要求 $t_{12}'=0$，从而
\begin{equation*}
  s_{12}^2=-l_{12}^{\prime 2}<0.
\end{equation*}

因之，仅当两个事件的间隔 $s_{12}$ 是虚数的情况下，我们才可以找到所要求的参考系.虚数间隔称为类空间隔.

因此，若两个事件的间隔是类空的，那么就有一个参考系存在，在其中，两个事件同时发生.在这个参考系中发生这两个事件的两点间的距离等于
\begin{equation}
  l_{12}'=\sqrt{l_{12}^2-c^2t_{12}^2}=\mathrm{i}s_{12}.
\end{equation}

由于间隔的不变性，将其分为类空间隔及类时间隔具有绝对的意义.这就是说，一个间隔的类空性或类时性与参考系无关.

取某一个事件 $O$ 作为时间及空间坐标的原点.换句话说，其轴上标有 $x,y,z,t$ 的四维坐标系中，事件 $O$ 的世界点就是坐标系的原点.现在我们来研究所有其他事件对于本事件 $O$ 的关系.为了便于作图说明，我们只考虑一维空间与时间，并将它们取在图2的两个轴上.一个当 $t=0$ 时经过 $x=0$ 的粒子的匀速直线运动，可以用一条直线来表示，这条直线经过 $O$ 点，与 $t$ 轴成一角，此角的正切等于粒子的速度.因为最大可能的速度是 $c$，所以这条直线与 $t$ 轴所成之角也有一个最大值.图2中有两条直线，代表两个信号经过事件 $O$（即当 $t=0$ 时经过 $x=0$）（以光速）向相反两个方向的传播.所有代表粒子运动的直线只能在 $aOc$ 及 $dOb$ 两个区域内.显然，在直线 $ab$ 及 $cd$ 上，$x=\pm ct$.我们先研究世界点在 $aOc$ 区域内的那些事件.很容易理解，在这个区域内的所有的点，$c^2t^2-x^2>0$.换句话说，在这个区域内的任何事件与 $O$ 事件的间隔是类时的.在这个区域内 $t>0$，即其中所有的事件都发生在 $O$ 事件之“后”.但是两个事件若被类时间隔所分开，无论在哪一个参考系内都不可能同时发生.可见，不可能找到一个参考系，其中 $aOc$ 区域的任何事件会在事件 $O$ 之“前”发生，即在 $t<0$ 时发生.因此，在 $aOc$ 区域内的所有事件对 $O$ 来说在任何参考系中都是未来的事件.所以，这个区域对事件 $O$ 来说可称为绝对未来.

\par\nopagebreak\noindent\begin{minipage}{\linewidth}\centering\includegraphics[width=5cm]{fig2.pdf}\\[0.3em]{\small 图 2}\end{minipage}\par\nopagebreak

完全类似地，所有在区域 $bOd$ 内的事件对 $O$ 来说都是绝对过去，即本区域内的事件，无论在任何参考系中都在 $O$ 事件前发生.

最后，再研究 $dOa$ 及 $cOb$ 两个区域.本区域内的任何事件与 $O$ 事件的间隔都是类空的.这些事件在任何参考系中发生在空间里的不同点.因此，这些区域对于 $O$ 来说都可称为绝对分隔.但是，关于这些事件的“同时”、“较早”及“较晚”等概念都是相对的.对这些区域内的任何事件来说，在一些参考系中，此事件在 $O$ 事件后发生；在另一些参考系中，此事件在 $O$ 事件前发生；最后也有一个参考系存在，在其中此事件与 $O$ 事件同时发生.

我们要注意，如果考虑所有三个空间坐标轴，而不只考虑一个空间坐标轴，那么，代替图2中的两条相交的线，在四维空间坐标系 $x,y,z,t$ 内，我们将有一个“圆锥体”$x^2+y^2+z^2-c^2t^2=0$，圆锥体的轴与 $t$ 轴重合（这个圆锥体称为光锥）.这时，绝对未来与绝对过去就由锥体的两个内区域来代表.

两个事件仅仅在其间隔是类时间隔的情况下，彼此才能有因果关系；这可由相互作用的传播速度不能大于光速这一事实直接推出来.如我们刚才已经看到的，正是对这些事件来说，“较早”与“较晚”等概念才有绝对意义，而这一点又是使因果概念具有意义的必要条件.

\section{固有时}

假设我们在某一个惯性参考系内观察一只钟，这只钟相对于我们可作任意形式的运动.在各个不同的时刻，该钟的运动可以认为是匀速的.因此，在每一时刻，我们可以引入一个固联于运动钟上的坐标系，这个坐标系（和该钟在一起）也是一个惯性参考系.

在无限小的时间间隔 $\mathrm{d}t$ 内（根据我们静止系内钟的读数），运动的钟所行进的距离是 $\sqrt{\mathrm{d}x^2+\mathrm{d}y^2+\mathrm{d}z^2}$.我们要问，这段时间运动的钟所指示的时间间隔 $\mathrm{d}t'$ 又如何？在与运动的钟固联的坐标系内，钟是静止的，即 $\mathrm{d}x'=\mathrm{d}y'=\mathrm{d}z'=0$，因为间隔是不变的，所以
\begin{equation*}
  \mathrm{d}s^2=c^2\mathrm{d}t^2-\mathrm{d}x^2-\mathrm{d}y^2-\mathrm{d}z^2=c^2\mathrm{d}t'^2,
\end{equation*}
由此可得
\begin{equation*}
  \mathrm{d}t'=\mathrm{d}t\sqrt{1-\frac{\mathrm{d}x^2+\mathrm{d}y^2+\mathrm{d}z^2}{c^2\mathrm{d}t^2}}.
\end{equation*}
但是
\begin{equation*}
  \frac{\mathrm{d}x^2+\mathrm{d}y^2+\mathrm{d}z^2}{\mathrm{d}t^2}=v^2,
\end{equation*}
其中 $v$ 为运动的钟的速度；所以
\begin{equation}
  \mathrm{d}t'=\frac{\mathrm{d}s}{c}=\mathrm{d}t\sqrt{1-\frac{v^2}{c^2}}.
\end{equation}

将上式积分，我们可以得到，当静止的钟所行走的时间为 $t_2-t_1$ 时，运动的钟所指示的时间间隔是
\begin{equation}
  t_2'-t_1'=\int_{t_1}^{t_2}\mathrm{d}t\sqrt{1-\frac{v^2}{c^2}}.
\end{equation}

随着某一给定物体一同运动的钟所指示的时间，称为该物体的固有时.公式（3.1）及（3.2）是将固有时通过观察运动的参考系的时间表示出来.

由（3.1）及（3.2）可见，一个运动物体的固有时永远较在静止系内相对应的时间间隔为小.换句话说，运动的钟较静止的钟走得慢些.

假设有一只钟相对于某一惯性参考系 $K$ 作匀速直线运动.一个同这只钟固联着的参考系 $K'$ 也是惯性系.从 $K$ 系内的观察者来看，$K'$ 系内的钟走慢了.反过来说，从 $K'$ 系内的观察者来看，$K$ 系内的钟走慢了.为了使我们相信这是不矛盾的，可以注意下面的事实.为了确定 $K'$ 系内的钟比 $K$ 系内的钟慢，我们必须按下述方法去做.假设在某一时刻，$K'$ 内的那只钟经过 $K$ 内的一只钟的旁边，而在这一时刻，两只钟所指示的时间恰好一样.为了比较 $K$ 及 $K'$ 内钟的快慢，我们必须再次将 $K'$ 内同一只动钟的读数与 $K$ 内的钟的读数作比较.但是在新的时刻，$K'$ 内的钟将从 $K$ 内的另一些钟旁边经过，现在我们就将运动的钟与那些钟比较.这时我们发现，$K'$ 内的钟比借以比较的 $K$ 内的钟走得慢.由此可见，为了比较两个参考系内的钟的快慢，我们需要在一个参考系内有几只钟，而在另一个参考系内有一只钟.因此这种过程对这两个参考系来说，并不是对称的.与另一参考系内不同的一些钟相比较的那一只钟总是走得慢的钟.

假设我们有两只钟，其中之一描绘一闭合路径，又回到出发点（即静止的钟所在之点），显然运动的钟慢了（与静止的钟相比）.相反，若设想运动的钟静止（而静止的钟运动），就不能做上面的推论了，因为那只钟既然描绘了一条闭合曲线，它的运动就不是匀速直线运动，因而与之相连的参考系就不是惯性系.

因为自然定律只有在不同的惯性系内才是一样的，与静止的钟相连的系统（惯性系）及与运动的钟相连的系统（非惯性系）具有不同的特性，因而导致静止的钟应当变慢这个结论的论证就不对了.

一只钟所指示的时间间隔，等于沿着钟的世界线而取的积分 $\frac{1}{c}\int\mathrm{d}s$.假如钟是静止的，则显然它的世界线是一条与 $t$ 轴平行的直线；假如钟在闭合路径上作非匀速运动而且又回到出发点，那么，它的世界线就是一条曲线，这条曲线经过静止钟的直的世界线的两点，这两点对应运动的起点及终点.另一方面，我们看到，静止钟所指示的时间间隔永远较运动钟所指示的为大.因此，我们得到一个结论，在两个世界点间所取的积分 $\int\mathrm{d}s$，如果是沿着连接这两点的直的世界线进行则有最大值\footnote{当然，需假设 $a$ 及 $b$ 两点和连接它们的曲线满足如下要求：沿该曲线的所有线元 $\mathrm{d}s$ 都是类时的.积分的这个性质同四维几何的伪欧几里得特性有关.在欧几里得空间中，这个积分沿直线当然取最小值.}.

\section{洛伦兹变换}

我们现在的目的是要找出从一个惯性系到另一个惯性系的变换公式，根据这个公式，当某一个事件在 $K$ 系内的坐标 $x,y,z,t$ 为已知时，就可以找到同一事件在另一个惯性系 $K'$ 内的坐标 $x',y',z',t'$.

在经典力学中，这个问题可以很简单地解决.由于时间的绝对性，我们有 $t=t'$；其次，假若坐标轴的取法和通常一样（即 $x,x'$ 两轴重合，$y,z$ 分别与 $y',z'$ 轴平行，运动是沿着 $x,x'$ 轴），那么很明显，$y,z$ 就等于 $y',z'$，而 $x$ 与 $x'$ 则相差一个距离，即一个坐标系相对于另一坐标系所走的距离.假如两个坐标系重合的时刻被取为时间的原点，并假设 $K'$ 系对于 $K$ 系的相对运动的速度为 $V$，那么，这个距离就是 $Vt$.所以，
\begin{equation}
  x=x'+Vt,\quad y=y',\quad z=z',\quad t=t'.
\end{equation}

这些公式称为伽利略变换.很容易证明，这个变换式，如我们所预料的一样，不能满足相对论的要求；这个变换式不能使事件与事件之间的间隔不变.

我们将从事件之间间隔不变的要求出发，推出相对论的变换公式.

正如在 §2 里所看到的，二事件间的间隔可以认为是在四维空间内的相对应的两个世界点间的距离.因此我们可以说，所要求的变换，必须使在四维空间 $x,y,z,ct$ 内的所有距离不变.但是这些变换仅仅包含坐标系的平移与转动.其中，我们对于坐标轴相对于自己作的平移并无兴趣，因为这不过是将空间坐标的原点移动一下，并将时间的参考点改变一下而已.所以，所要求的变换，在数学上应当表示为四维坐标系 $x,y,z,ct$ 的转动.

四维空间内的一切转动可以分解为六个分别在六个平面 $xy,zy,xz,tx,ty,tz$ 内的转动（正如在三维空间内的一切转动可以分解为三个分别在 $xy,yz,xz$ 三个平面内的转动一样）.其中，前三个转动仅仅变换空间坐标，它们对应通常的空间转动.

我们研究在 $tx$ 平面内的转动，这时 $y$ 与 $z$ 坐标是不变的.具体地说，这个变换必须使差值 $(ct)^2-x^2$，即点 $(ct,x)$ 到原点“距离”的平方保持不变.新旧坐标的关系最一般地由以下二式决定：
\begin{equation}
  x=x'\cosh\psi+ct'\sinh\psi,\quad
  ct=x'\sinh\psi+ct'\cosh\psi,
\end{equation}
式中 $\psi$ 为转动角；简单验算表明，实际上有 $c^2t^2-x^2=c^2t'^2-x'^2$.(4.2) 式与坐标轴转动变换的通常公式不同之处在于，后者中的三角函数换成了双曲函数.这就是伪欧几里得几何与欧几里得几何的差别.

我们现在要找出由一个惯性参考系 $K$ 到另外一个惯性系 $K'$ 的变换公式，$K'$ 以速度 $V$ 沿 $x$ 轴对 $K$ 作相对运动.在此情况下，显然只有空间坐标 $x$ 与时间 $t$ 发生变化.所以这个变换必须有（4.2）的形式.现在只剩下决定转动角 $\psi$ 的问题，$\psi$ 仅与相对速度 $V$ 有关\footnote{注意，为了避免弄混，以后永远用 $V$ 表示两个惯性系相对运动的恒定速度，用 $v$ 表示粒子运动的速度，$v$ 并不必须为常数.}.

我们来研究参考系 $K'$ 的原点在 $K$ 内的运动.这时 $x'=0$，而公式（4.2）可写成
\begin{equation*}
  x=ct'\sinh\psi,\quad ct=ct'\cosh\psi,
\end{equation*}
相除可得
\begin{equation*}
  \frac{x}{ct}=\tanh\psi.
\end{equation*}
但 $x/t$ 显然是 $K'$ 对 $K$ 的速度 $V$.因此，
\begin{equation*}
  \tanh\psi=\frac{V}{c}.
\end{equation*}
由此得
\begin{equation*}
  \sinh\psi=\frac{\dfrac{V}{c}}{\sqrt{1-\dfrac{V^2}{c^2}}},\quad
  \cosh\psi=\frac{1}{\sqrt{1-\dfrac{V^2}{c^2}}}.
\end{equation*}
代入（4.2），得：
\begin{equation}
  x=\frac{x'+Vt'}{\sqrt{1-\dfrac{V^2}{c^2}}},\quad
  y=y',\quad z=z',\quad
  t=\frac{t'+\dfrac{V}{c^2}x'}{\sqrt{1-\dfrac{V^2}{c^2}}}.
\end{equation}

这就是所要求的变换公式.它们被称为洛伦兹变换，是今后讨论的基础.

用 $x,y,z,t$ 来表示 $x',y',z',t'$ 的逆公式只需以 $-V$ 代替 $V$ 便得（因为 $K$ 系以速度 $-V$ 相对 $K'$ 运动）.这些公式也可直接求解方程式（4.3）得到.

由（4.3）式易见，取 $c\to\infty$ 的经典力学极限，洛伦兹变换事实上就过渡到伽利略变换了.

当 $V>c$ 时，（4.3）式中的 $x,t$ 变成虚数；这与运动速度不可能大于光速的事实符合.此外，我们也不可以用以光速运动的参考系，因为在这种情形下，（4.3）式的分母将为零.

当 $V$ 比光速小很多时，我们可以用下面的近似公式代替（4.3）：
\begin{equation}
  x=x'+Vt',\quad y=y',\quad z=z',\quad t=t'+\frac{V}{c^2}x'.
\end{equation}

假设在 $K$ 系内有一根平行于 $x$ 轴的静止杆.假定它在 $K$ 系内测定的长度为 $\Delta x=x_2-x_1$（$x_2$ 及 $x_1$ 为杆两端在 $K$ 系内的坐标）.我们现在来求此杆在 $K'$ 系内的长度.为此目的，我们需要在同一时刻 $t'$ 找出杆两端在 $K'$ 内的坐标 $x_2'$ 及 $x_1'$.由（4.3）我们得到
\begin{equation*}
  x_1=\frac{x_1'+Vt'}{\sqrt{1-\dfrac{V^2}{c^2}}},\quad
  x_2=\frac{x_2'+Vt'}{\sqrt{1-\dfrac{V^2}{c^2}}}.
\end{equation*}
杆在 $K'$ 内的长度是 $\Delta x'=x_2'-x_1'$；由 $x_2$ 减去 $x_1$，得
\begin{equation*}
  \Delta x=\frac{\Delta x'}{\sqrt{1-\dfrac{V^2}{c^2}}}.
\end{equation*}
杆的固有长度是它在相对它静止的参考系内的长度.以 $l_0=\Delta x$ 代表这个固有长度，以 $l$ 代表它在任何其他参考系 $K'$ 内的长度.那么，
\begin{equation}
  l=l_0\sqrt{1-\frac{V^2}{c^2}}.
\end{equation}

因此，一根杆在它是静止的那个参考系内最长.在它是以速度 $V$ 运动的那个参考系内，它的长度就要减少一个因子 $\sqrt{1-\dfrac{V^2}{c^2}}$.相对论的这个结果称为洛伦兹收缩.

因为物体的横向尺度（如宽及高）都不因运动而变，所以它的体积 $\mathcal{V}$ 也按照相似的公式收缩，即
\begin{equation}
  \mathcal{V}=\mathcal{V}_0\sqrt{1-\frac{V^2}{c^2}},
\end{equation}
其中 $\mathcal{V}_0$ 代表物体的固有体积.

由洛伦兹变换，还可以得到我们已经知道的有关固有时（§3）的结果.假设在 $K'$ 系内有一只静止的钟.我们假定有两个事件发生在 $K'$ 内同一空间点 $x',y',z'$.在 $K'$ 内这两事件之间的时间为 $\Delta t'=t_2'-t_1'$.现在我们要找出，在 $K$ 系内，同样的两事件之间的时间 $\Delta t$.由（4.3）得
\begin{equation*}
  t_1=\frac{t_1'+\dfrac{V}{c^2}x'}{\sqrt{1-\dfrac{V^2}{c^2}}},\quad
  t_2=\frac{t_2'+\dfrac{V}{c^2}x'}{\sqrt{1-\dfrac{V^2}{c^2}}},
\end{equation*}
相减则得
\begin{equation*}
  t_2-t_1=\Delta t=\frac{\Delta t'}{\sqrt{1-\dfrac{V^2}{c^2}}},
\end{equation*}
这一公式与公式（3.1）完全符合.

最后，我们再谈谈洛伦兹变换有别于伽利略变换的另一个一般性质.后者具有可对易性，即接连两次伽利略变换（具有不同速度 $V_1$ 和 $V_2$）的联合结果与施行变换的顺序无关.另一方面，接连两次洛伦兹变换的结果一般却依赖于它们的顺序.这一点已经可以从我们将这些变换公式描述为四维坐标系的转动而纯数学地看出来：我们知道，两次转动（绕不同轴）的结果依赖于施行它们的顺序.唯一的例外是矢量 $V_1$ 和 $V_2$ 平行的变换情况（这等价于四维坐标系绕相同轴的两次转动）.

\section{速度的变换}

在上节中，我们求得一些公式，用这些公式，我们能够从一个事件在一个参考系内的坐标找出同一事件在第二个参考系内的坐标.现在我们要求出公式，用以表示一个实物粒子在一个参考系内的速度与其在第二个参考系内的速度之间的关系.

我们再一次假定 $K'$ 系相对于 $K$ 系以速度 $V$ 沿 $x$ 轴运动.设 $v_x=\mathrm{d}x/\mathrm{d}t$ 为粒子在系统 $K$ 内的速度的分量，而 $v_x'=\mathrm{d}x'/\mathrm{d}t'$ 为同一粒子在系统 $K'$ 内的速度的分量.由（4.3），得
\begin{equation*}
  \mathrm{d}x=\frac{\mathrm{d}x'+V\mathrm{d}t'}{\sqrt{1-\dfrac{V^2}{c^2}}},\quad
  \mathrm{d}y=\mathrm{d}y',\quad \mathrm{d}z=\mathrm{d}z',\quad
  \mathrm{d}t=\frac{\mathrm{d}t'+\dfrac{V}{c^2}\mathrm{d}x'}{\sqrt{1-\dfrac{V^2}{c^2}}}.
\end{equation*}
用第四个方程除前三个并引入速度
\begin{equation*}
  \boldsymbol{v}=\frac{\mathrm{d}\boldsymbol{r}}{\mathrm{d}t},\quad
  \boldsymbol{v}'=\frac{\mathrm{d}\boldsymbol{r}'}{\mathrm{d}t'}
\end{equation*}
则得
\begin{equation}
  v_x=\frac{v_x'+V}{1+v_x'\dfrac{V}{c^2}},\quad
  v_y=\frac{v_y'\sqrt{1-\dfrac{V^2}{c^2}}}{1+v_x'\dfrac{V}{c^2}},\quad
  v_z=\frac{v_z'\sqrt{1-\dfrac{V^2}{c^2}}}{1+v_x'\dfrac{V}{c^2}}.
\end{equation}

这些公式就决定了速度的变换.它们是相对论里的速度合成法则.在极限情形下，即 $c\to\infty$ 时，它们就变为经典力学里的公式：
\begin{equation*}
  v_x=v_x'+V,\quad v_y=v_y',\quad v_z=v_z'.
\end{equation*}

在粒子沿 $x$ 轴运动的特殊情况下，$v_x=v$，$v_y=v_z=0$.那么，$v_y'=v_z'=0$，$v_x'=v'$，并且
\begin{equation}
  v=\frac{v'+V}{1+v'\dfrac{V}{c^2}}.
\end{equation}

很容易证明，如果两个速度各小于或等于光速，其合成速度，根据这个公式，也不会大于光速.

假如速度 $V$ 比光速 $c$ 小很多（$v$ 可以是任意的），我们将近似地（精确到 $V/c$ 的项）得到
\begin{equation*}
  v_x=v_x'+V\left(1-\frac{v_x'^2}{c^2}\right),\quad
  v_y=v_y'-v_x'v_y'\frac{V}{c^2},\quad
  v_z=v_z'-v_x'v_z'\frac{V}{c^2}.
\end{equation*}
这3个公式可以简写成一个矢量公式
\begin{equation}
  \boldsymbol{v}=\boldsymbol{v}'+\boldsymbol{V}
  -\frac{1}{c^2}(\boldsymbol{V}\cdot\boldsymbol{v}')\boldsymbol{v}'.
\end{equation}

我们可以指出，在相对论的速度合成公式（5.1）中，相加的两个速度 $\boldsymbol{v}'$ 和 $\boldsymbol{V}$ 是以不对称的方式引入的（倘若它们不都沿 $x$ 轴指向的话）.这个事实同洛伦兹变换的非对易性有关，我们将在下节提到.

让我们这样来选择坐标轴，使粒子的速度在给定时刻是在 $xy$ 平面内.这时粒子在 $K$ 系内的速度分量是 $v_x=v\cos\theta$，$v_y=v\sin\theta$，而在 $K'$ 内则为 $v_x'=v'\cos\theta'$，$v_y'=v'\sin\theta'$（$v,v'$ 为速度在 $K$ 及 $K'$ 内的绝对值；$\theta,\theta'$ 为速度与 $x$ 轴及 $x'$ 轴所夹之角）.用（5.1）式，我们就得到
\begin{equation}
  \tan\theta=\frac{v'\sqrt{1-\dfrac{V^2}{c^2}}\sin\theta'}
  {v'\cos\theta'+V}.
\end{equation}

这个公式决定了速度的方向从一个参考系变换到另一个参考系时的改变.

让我们来详尽地研究这个公式的一个重要特例，即光由一个参考系变换到另一个参考系时的偏差，即所谓光行差现象.在这种情形下 $v=v'=c$，因而（5.4）式化为
\begin{equation}
  \tan\theta=\frac{\sqrt{1-\dfrac{V^2}{c^2}}}{\dfrac{V}{c}+\cos\theta'}\sin\theta'.
\end{equation}

由同一变换公式（5.1），用相似的方法，很容易得到
\begin{equation}
  \sin\theta=\frac{\sqrt{1-\dfrac{V^2}{c^2}}}{1+\dfrac{V}{c}\cos\theta'}\sin\theta',\quad
  \cos\theta=\frac{\cos\theta'+\dfrac{V}{c}}{1+\dfrac{V}{c}\cos\theta'}.
\end{equation}

如果 $V\ll c$，由（5.6）式我们可以得到精确到数量级为 $V/c$ 项的公式如下：
\begin{equation*}
  \sin\theta-\sin\theta'=-\frac{V}{c}\sin\theta'\cos\theta'.
\end{equation*}
若引入 $\Delta\theta=\theta'-\theta$（光行差角），我们就得到同级的近似公式
\begin{equation}
  \Delta\theta=\frac{V}{c}\sin\theta',
\end{equation}
这就是著名的光行差的基本公式.

\section{四维矢量}

一个事件的坐标 $(ct,x,y,z)$ 可以看成四维空间中一个四维径向矢量的分量.我们将把它的分量记为 $x^i$，这里指标 $i$ 取值 $0,1,2,3$，而且
\begin{equation*}
  x^0=ct,\quad x^1=x,\quad x^2=y,\quad x^3=z.
\end{equation*}
该径向四维矢量“长度”的平方由下式给出
\begin{equation*}
  (x^0)^2-(x^1)^2-(x^2)^2-(x^3)^2.
\end{equation*}
它在四维坐标系的任意转动下不变，特别是，它在洛伦兹变换下不变.

更一般地，如果4个量 $A^0,A^1,A^2,A^3$，在四维坐标系的变换下像四维径向矢量的分量 $x^i$ 那样变换，我们就将这4个量的集合称为四维矢量 $A^i$.在洛伦兹变换下，
\begin{equation}
  A^0=\frac{A^{\prime 0}+\dfrac{V}{c}A^{\prime 1}}{\sqrt{1-\dfrac{V^2}{c^2}}},\quad
  A^1=\frac{A^{\prime 1}+\dfrac{V}{c}A^{\prime 0}}{\sqrt{1-\dfrac{V^2}{c^2}}},\quad
  A^2=A^{\prime 2},\quad A^3=A^{\prime 3}.
\end{equation}

与四维径向矢量的平方类似，任一四维矢量数值的平方定义为：
\begin{equation*}
  (A^0)^2-(A^1)^2-(A^2)^2-(A^3)^2.
\end{equation*}

为表示方便起见，我们引入四维矢量分量的两种“类型”，用带上标和下标的符号 $A^i$ 和 $A_i$ 来标记它们.两者之间的关系是
\begin{equation}
  A_0=A^0,\quad A_1=-A^1,\quad A_2=-A^2,\quad A_3=-A^3.
\end{equation}

量 $A^i$ 称为四维矢量的逆变分量，$A_i$ 称为协变分量.四维矢量的平方则取形式
\begin{equation*}
  \sum_{i=0}^{3}A^iA_i=A^0A_0+A^1A_1+A^2A_2+A^3A_3.
\end{equation*}

人们通常略去求和号，将这样的求和简单记为 $A^iA_i$.也就是约定遍历所有重复指标求和，而把求和号省去.每对指标中必须一个为上标，另一个为下标.这种遍历“傀”指标求和的约定非常方便，可大大简化公式的书写.

我们将用拉丁字母 $i,k,l,\cdots$ 表示四维指标，取值 $0,1,2,3$.

与四维矢量的平方类比，我们可以构造两个不同四维矢量的标积：
\begin{equation*}
  A^iB_i=A^0B_0+A^1B_1+A^2B_2+A^3B_3.
\end{equation*}
显然，这既可以写为 $A^iB_i$，也可以写为 $A_iB^i$，结果相同.我们一般可以交换任何一对傀指标中的上标和下标\footnote{现代文献中常常省略四维矢量的指标，将它们的平方与标积写为 $A^2$ 和 $AB$.在本教程中我们将不用这些符号.}.

积 $A^iB_i$ 是一个四维标量——它在四维坐标系的转动下是不变的.这一点很容易直接验证\footnote{应当记得，以协变分量表示的四维矢量变换法则（在正负号上）不同于以逆变分量表示的同一法则.因此，代替（6.1），我们有：
\begin{equation*}
  A_0=\frac{A'_0-\dfrac{V}{c}A'_1}{\sqrt{1-\dfrac{V^2}{c^2}}},\quad
  A_1=\frac{A'_1-\dfrac{V}{c}A'_0}{\sqrt{1-\dfrac{V^2}{c^2}}},\quad
  A_2=A'_2,\quad A_3=A'_3.
\end{equation*}}，但也可从所有四维矢量都按相同法则变换的事实预先看出（与平方 $A^iA_i$ 类比）.

分量 $A^0$ 称为四维矢量的时间分量，$A^1,A^2,A^3$ 称为四维矢量的空间分量（与四维径向矢量类比）.四维矢量的平方可以为正、负或零；这样的矢量分别称为类时矢量、类空矢量和类光矢量（类似于对间隔所用的术语）\footnote{类光矢量也称各向同性矢量.}.

在纯空间转动（即不影响时间轴的变换）下，四维矢量 $A^i$ 的三个空间分量构成一个三维矢量 $\boldsymbol{A}$.该四维矢量的时间分量（在这些变换下）是一个三维标量.为了列举四维矢量的分量，我们常将其写为
\begin{equation*}
  A^i=(A^0,\boldsymbol{A}).
\end{equation*}
同一四维矢量的协变分量为 $A_i=(A^0,-\boldsymbol{A})$.该四维矢量的平方是 $A^iA_i=(A^0)^2-\boldsymbol{A}^2$.因此，对于四维径向矢量：
\begin{equation*}
  x^i=(ct,\boldsymbol{r}),\quad x_i=(ct,-\boldsymbol{r}),\quad
  x^ix_i=c^2t^2-\boldsymbol{r}^2.
\end{equation*}

对于三维矢量（带坐标 $x,y,z$），没有必要区分逆变和协变分量.只要能够做到这一点而不致引起混淆，我们将用希腊字母作为下标把这些分量记为 $A_\alpha\ (\alpha=x,y,z)$.特别是我们将假设对于任何重复指标遍历 $x,y,z$ 求和（例如，$\boldsymbol{A}\cdot\boldsymbol{B}=A_\alpha B_\alpha$）.

二阶四维张量是16个量 $A^{ik}$ 的集合，它在坐标变换下像两个四维矢量分量的积那样变换.我们可以类似地定义更高阶的四维张量.

一个二阶张量的分量可以写为三种形式：协变的 $A_{ik}$，逆变的 $A^{ik}$ 和混合的 $A^i{}_k$（这里，在后一种情形下，应当区分 $A^i{}_k$ 和 $A_i{}^k$，即应当仔细搞清楚两个指标中哪一个是上标，哪一个是下标）.不同类型分量之间的联系由以下通则决定：升或降一个空间指标（1,2,3）改变分量的正负号，而升或降时间指标（0）则不变号.因此：
\begin{equation*}
\begin{gathered}
  A_{00}=A^{00},\quad A_{01}=-A^{01},\ A_{11}=A^{11},\cdots,\\
  A^0{}_0=A^{00},\quad A_0{}^1=A^{01},\quad A^0{}_1=-A^{01},\quad
  A^1{}_1=-A^{11},\cdots
\end{gathered}
\end{equation*}

在纯空间变换下，9个量 $A^{11},A^{12},\cdots$ 构成一个三维张量.三个分量 $A^{01},A^{02},A^{03}$ 和三个分量 $A^{10},A^{20},A^{30}$ 构成三维矢量，而分量 $A^{00}$ 是一个三维标量.

如果 $A^{ik}=A^{ki}$，张量 $A^{ik}$ 称为对称的；如果 $A^{ik}=-A^{ki}$，则称反对称的.在反对称张量中，所有对角分量（即分量 $A^{00},A^{11},\cdots$）都是零，因为，例如，我们必须有 $A^{00}=-A^{00}$.对于一个对称张量 $A^{ik}$，混合分量 $A^i{}_k$ 和 $A_i{}^k$ 显然重合；在这样的情形下，我们把一个指标置于另一个上方，简单地记为 $A^i{}_k$.

在每个张量方程中，等号两边所含的自由指标（区别于傀指标）必须字母相同且位置相同（即上或下）.张量方程中的自由指标可以上移或下移，但必须对方程中所有的项同时进行.让不同张量的协变和逆变分量相等是“非法的”；这样的方程即便碰巧在特定参考系中成立，在变换到另一个参考系时也会失效.

通过对张量 $A^{ik}$ 的分量求和可以形成一个标量
\begin{equation*}
  A^i{}_i=A^0{}_0+A^1{}_1+A^2{}_2+A^3{}_3
\end{equation*}
（当然，这里 $A^i{}_i=A_i{}^i$）.这个和称为张量的迹，求得它的运算称为缩并.

前面讨论的两个四维矢量的标积的形成就是一个缩并运算：它是从张量 $A^iB_k$ 形成标量 $A^iB_i$.缩并任何一对指标一般会使张量的阶减去2.例如，$A^i{}_{kli}$ 是一个二阶张量，$A^i{}_kB^k$ 是一个四维矢量，$A^{ik}{}_{ik}$ 是一个标量，等等.

单位四维张量 $\delta^i_k$ 满足如下条件：对于任意四维矢量 $A^i$
\begin{equation}
  \delta^k_iA^i=A^k.
\end{equation}
这个张量的分量显然是
\begin{equation}
  \delta^k_i=
  \begin{cases}
    1, & \text{当 } i=k,\\
    0, & \text{当 } i\neq k.
  \end{cases}
\end{equation}
它的迹是 $\delta^i_i=4$.

通过在 $\delta^k_i$ 中升一个指标或降另一个指标，我们可以得到逆变张量 $g^{ik}$ 或协变张量 $g_{ik}$，称之为度规张量.张量 $g^{ik}$ 和 $g_{ik}$ 具有相同的分量，可以写成矩阵：
\begin{equation}
  (g^{ik})=(g_{ik})=
  \begin{bmatrix}
    1 & 0 & 0 & 0\\
    0 & -1 & 0 & 0\\
    0 & 0 & -1 & 0\\
    0 & 0 & 0 & -1
  \end{bmatrix}
\end{equation}
（指标 $i$ 标记行，$k$ 标记列，顺序为 $0,1,2,3$）.显然有
\begin{equation}
  g_{ik}A^k=A_i,\quad g^{ik}A_k=A^i.
\end{equation}

两个四维矢量的标积因而可以写成形式：
\begin{equation}
  A^iA_i=g_{ik}A^iA^k=g^{ik}A_iA_k.
\end{equation}

张量 $\delta^i_k,g_{ik}$ 和 $g^{ik}$ 的特别之处在于，它们的分量在所有坐标系中都相同.四阶的全反对称单位张量 $e^{iklm}$ 具有同样性质.这个张量的分量在交换任一对指标时变号，其非零分量为 $\pm 1$.从反对称性可知，有两个指标相同的所有分量均为零，所以，仅有的非零分量是那些所有4个指标都不同者.我们令
\begin{equation}
  e^{0123}=+1
\end{equation}
（所以，$e_{0123}=-1$）.于是，所有其他的非零分量 $e^{iklm}$ 等于 $+1$ 或 $-1$，依 $i,k,l,m$ 这几个数能经偶数还是奇数次换位排成 $0,1,2,3$ 而定.这样的分量数是 $4!=24$，所以，
\begin{equation}
  e^{iklm}e_{iklm}=-24.
\end{equation}

对于坐标系的转动而言，$e^{iklm}$ 诸量的特性与张量分量的特性相同，但是如果我们改变1个或3个坐标的正负号，分量 $e^{iklm}$ 并不改变，因为按定义它们在所有坐标系中都相同，而张量的分量在这种情况下是应当变号的.所以，严格地说，$e^{iklm}$ 并不是张量，而是一个赝张量.任意阶的赝张量，特别是赝标量，在所有的坐标变换下都具有张量的性质，只有那些不能归结为转动的变换，即反射（不能归结为转动的坐标正负号改变）是例外.

乘积 $e^{iklm}e^{prst}$ 构成一个8阶四维张量，它是一个真正的张量.通过缩并一对或多对指标可以得到6阶、4阶和2阶张量.所有这些张量在所有坐标系中具有相同的形式.所以，它们的分量必须表示为单位张量 $\delta^i_k$（其分量在所有坐标系中都相同的唯一真张量）分量乘积的组合.从指标排列必须具有的对称性出发，这些组合是不难求得的.\footnote{我们给出下列公式以备参考：
\begin{equation*}
  e^{iklm}e_{prst}=-\begin{vmatrix}
    \delta^i_p & \delta^i_r & \delta^i_s & \delta^i_t\\
    \delta^k_p & \delta^k_r & \delta^k_s & \delta^k_t\\
    \delta^l_p & \delta^l_r & \delta^l_s & \delta^l_t\\
    \delta^m_p & \delta^m_r & \delta^m_s & \delta^m_t
  \end{vmatrix},\qquad
  e^{iklm}e_{prsm}=-\begin{vmatrix}
    \delta^i_p & \delta^i_r & \delta^i_s\\
    \delta^k_p & \delta^k_r & \delta^k_s\\
    \delta^l_p & \delta^l_r & \delta^l_s
  \end{vmatrix},
\end{equation*}
\begin{equation*}
  e^{iklm}e_{prlm}=-2(\delta^i_p\delta^k_r-\delta^i_r\delta^k_p),\qquad
  e^{iklm}e_{pklm}=-6\delta^i_p.
\end{equation*}
这些公式中的整体系数可以用（6.9）式给出的完全缩并的结果来核对.
由这些公式我们得出如下结果：
\begin{equation*}
  e^{prst}A_{ip}A_{kr}A_{ls}A_{mt}=-Ae_{iklm},\qquad
  e^{iklm}e^{prst}A_{ip}A_{kr}A_{ls}A_{mt}=24A,
\end{equation*}
式中 $A$ 是由量 $A_{ik}$ 形成的行列式.}

如果 $A^{ik}$ 是一个反对称张量，则张量 $A^{ik}$ 和赝张量 $A^{*ik}=(1/2)e^{iklm}A_{lm}$ 称为彼此对偶.类似地，$e^{iklm}A_m$ 是一个与矢量 $A^i$ 对偶的三阶反对称赝张量.对偶张量的乘积 $A^{ik}A^{*}_{ik}$ 显然是一个赝标量.

联系这里的讨论，我们来提一提三维矢量和张量的一些类似性质.三阶全反对称单位赝张量 $e_{\alpha\beta\gamma}$ 是这样一些量的集合，它们在任何一对指标换位时变号.在 $e_{\alpha\beta\gamma}$ 的分量中，只有那些具有3个不同指标者才不等于零.我们令 $e_{xyz}=1$；其他的分量等于 $1$ 或 $-1$ 则依 $\alpha,\beta,\gamma$ 这个序列能经偶数还是奇数次换位排成 $x,y,z$ 的顺序而定.\footnote{四维张量 $e^{iklm}$ 的分量在四维坐标系的转动下不变，以及三维张量 $e_{\alpha\beta\gamma}$ 的分量在空间轴的转动下不变的事实，是如下普遍法则的特殊情况：任何阶数等于其定义空间维数的全反对称张量在该空间中坐标系的转动下是不变的.}

乘积 $e_{\alpha\beta\gamma}e_{\lambda\mu\nu}$ 构成一个6阶真三维张量，因而可以表示为单位三维张量 $\delta_{\alpha\beta}$ 分量乘积的组合.\footnote{我们给出下面的公式以备参考：
\begin{equation*}
  e_{\alpha\beta\gamma}e_{\lambda\mu\nu}=\begin{vmatrix}
    \delta_{\alpha\lambda} & \delta_{\alpha\mu} & \delta_{\alpha\nu}\\
    \delta_{\beta\lambda} & \delta_{\beta\mu} & \delta_{\beta\nu}\\
    \delta_{\gamma\lambda} & \delta_{\gamma\mu} & \delta_{\gamma\nu}
  \end{vmatrix}.
\end{equation*}
通过缩并该张量的一对、两对和三对指标，我们得到
\begin{equation*}
  e_{\alpha\beta\gamma}e_{\lambda\mu\gamma}=\delta_{\alpha\lambda}\delta_{\beta\mu}-\delta_{\alpha\mu}\delta_{\beta\lambda},\qquad
  e_{\alpha\beta\gamma}e_{\lambda\beta\gamma}=2\delta_{\alpha\lambda},\qquad
  e_{\alpha\beta\gamma}e_{\alpha\beta\gamma}=6.
\end{equation*}}

在坐标系的反射（即所有坐标变号）下，一个普通矢量也变号.这样的矢量称为极矢量.一个矢量若能写成两个极矢量的矢积，则其分量在反演下不变号.这样的矢量称为轴矢量.一个极矢量和一个轴矢量的标积并不是一个真标量，而是一个赝标量；它在坐标反演下变号.轴矢量是赝矢量，对偶于某反对称张量.因此，如果 $\boldsymbol{C}=\boldsymbol{A}\times\boldsymbol{B}$，那么
\begin{equation*}
  C_\alpha=\frac{1}{2}e_{\alpha\beta\gamma}C_{\beta\gamma},\qquad
  \text{其中}\quad C_{\beta\gamma}=A_\beta B_\gamma-A_\gamma B_\beta.
\end{equation*}

现在考虑四维张量.反对称张量 $A^{ik}$ 的空间分量（$i,k,\cdots=1,2,3$）对于纯空间变换构成一个三维反对称张量；根据我们的论述，其分量可以用一个三维轴矢量的分量来表示.对于同样的变换，分量 $A^{01},A^{02},A^{03}$ 构成一个三维极矢量.因此一个反对称四维张量可以写成矩阵：
\begin{equation}
  (A^{ik})=\begin{bmatrix}
    0 & p_x & p_y & p_z\\
    -p_x & 0 & -a_z & a_y\\
    -p_y & a_z & 0 & -a_x\\
    -p_z & -a_y & a_x & 0
  \end{bmatrix},
\end{equation}
这里，对于空间变换，$\boldsymbol{p}$ 和 $\boldsymbol{a}$ 分别为极矢量和轴矢量.在列出反对称四维张量的分量时，我们将它们写成形式
\begin{equation*}
  A^{ik}=(\boldsymbol{p},\boldsymbol{a});
\end{equation*}
同一张量的协变分量则为
\begin{equation*}
  A_{ik}=(-\boldsymbol{p},\boldsymbol{a}).
\end{equation*}

最后，我们来考虑四维张量分析的某些微分和积分运算.

标量 $\varphi$ 的四维梯度是四维矢量
\begin{equation*}
  \frac{\partial\varphi}{\partial x^i}
  =\left(\frac{1}{c}\frac{\partial\varphi}{\partial t},\nabla\varphi\right).
\end{equation*}
我们必须记住，这些导数是被看做该四维矢量的协变分量.事实上，该标量的微分
\begin{equation*}
  \mathrm{d}\varphi=\frac{\partial\varphi}{\partial x^i}\mathrm{d}x^i
\end{equation*}
也是一个标量.从它的形式（两个四维矢量的标积）看，我们的论断是显然的.

一般说来，对于坐标 $x^i$ 微分的算符 $\partial/\partial x^i$ 应当看成是该算符四维矢量的协变分量.因此，例如说，一个四维矢量的散度是一个标量，表达式为 $\partial A^i/\partial x^i$，其中我们是对逆变分量 $A^i$ 进行微分的.\footnote{如果我们对于“协变坐标”$x_i$ 进行微分，则导数
\begin{equation*}
  \frac{\partial\varphi}{\partial x_i}=g^{ik}\frac{\partial\varphi}{\partial x^k}
  =\left(\frac{1}{c}\frac{\partial\varphi}{\partial t},-\nabla\varphi\right)
\end{equation*}
构成一个四维矢量的逆变分量.我们将只在特殊情形下采用这种形式（例如，为了写出四维梯度的平方 $\frac{\partial\varphi}{\partial x^i}\frac{\partial\varphi}{\partial x_i}$）.
注意，在文献中对于坐标的偏导数常用符号缩写
\begin{equation*}
  \partial^i=\frac{\partial}{\partial x_i},\qquad
  \partial_i=\frac{\partial}{\partial x^i}.
\end{equation*}
以这种形式书写微分算符时，由它们构成的量的协变或逆变性质是一目了然的.另一种书写导数的缩写符号也有同样的优点，即在指标前面放一逗号：
\begin{equation*}
  \varphi_{,i}=\frac{\partial\varphi}{\partial x^i},\qquad
  \varphi^{,i}=\frac{\partial\varphi}{\partial x_i}.
\end{equation*}}

在三维空间中，可以沿体积、曲面或曲线进行积分.在四维空间中积分有四种类型：

1. 沿四维空间中一条曲线的积分.积分元就是线元，即四维矢量 $\mathrm{d}x^i$.

2. 沿四维空间中一个（二维）曲面的积分.我们知道，在三维空间中，由两个矢量 $\mathrm{d}\boldsymbol{r}$ 和 $\mathrm{d}\boldsymbol{r}'$ 构成的平行四边形的面积在坐标平面 $x_\alpha x_\beta$ 上的投影是 $\mathrm{d}x_\alpha\mathrm{d}x_\beta'-\mathrm{d}x_\beta\mathrm{d}x_\alpha'$.类似地，在四维空间中，无限小面元由二阶反对称张量 $\mathrm{d}f^{ik}=\mathrm{d}x^i\mathrm{d}x^{\prime k}-\mathrm{d}x^k\mathrm{d}x^{\prime i}$ 给定，其分量为该面元在坐标平面上的投影.众所周知，在三维空间中，人们不用张量 $\mathrm{d}f_{\alpha\beta}$ 而用与之对偶的矢量 $\mathrm{d}f_\alpha$ 来表示面元：$\mathrm{d}f_\alpha=\frac{1}{2}e_{\alpha\beta\gamma}\mathrm{d}f_{\beta\gamma}$.在几何上，这是一个与面元垂直的矢量，其绝对值等于面元的面积.在四维空间中，我们不能构造这样的矢量，但可以构造与张量 $\mathrm{d}f^{ik}$ 对偶的张量 $\mathrm{d}f^{*ik}$
\begin{equation}
  \mathrm{d}f^{*ik}=\frac{1}{2}e^{iklm}\mathrm{d}f_{lm}.
\end{equation}
在几何上，它描述这样一个面元，这个面元等于并“垂直”于面元 $\mathrm{d}f^{ik}$，其内的所有线段均正交于面元 $\mathrm{d}f^{ik}$ 内的所有线段.显然有 $\mathrm{d}f^{ik}\mathrm{d}f^{*}_{ik}=0$.

3. 沿一个超曲面，即沿一个三维流形的积分.在三维空间中，由3个矢量张成的平行六面体的体积等于由这些矢量的分量构成的3阶行列式.类似地，可以得到由3个四维矢量 $\mathrm{d}x^i,\mathrm{d}x^{\prime i},\mathrm{d}x^{\prime\prime i}$ 张成的平行六面体体积（即该超曲面的“面积”）的投影；它们由如下行列式给出
\begin{equation*}
  \mathrm{d}S^{ikl}=\begin{vmatrix}
    \mathrm{d}x^i & \mathrm{d}x^{\prime i} & \mathrm{d}x^{\prime\prime i}\\
    \mathrm{d}x^k & \mathrm{d}x^{\prime k} & \mathrm{d}x^{\prime\prime k}\\
    \mathrm{d}x^l & \mathrm{d}x^{\prime l} & \mathrm{d}x^{\prime\prime l}
  \end{vmatrix},
\end{equation*}
这构成一个3阶张量，对所有3个指标都是反对称的.像沿超曲面的积分元一样，使用与张量 $\mathrm{d}S^{ikl}$ 对偶的四维矢量 $\mathrm{d}S^i$ 更方便：
\begin{equation}
  \mathrm{d}S^i=-\frac{1}{6}e^{iklm}\mathrm{d}S_{klm},\quad
  \mathrm{d}S_{klm}=e_{nklm}\mathrm{d}S^n.
\end{equation}
这里
\begin{equation*}
  \mathrm{d}S^0=\mathrm{d}S^{123},\quad \mathrm{d}S^1=\mathrm{d}S^{023},\cdots
\end{equation*}
在几何上，$\mathrm{d}S^i$ 是一个四维矢量，数值上等于超曲面元的“面积”，并与该面元垂直（即垂直于该超曲面元内的所有线段）.特别是，$\mathrm{d}S^0=\mathrm{d}x\,\mathrm{d}y\,\mathrm{d}z$，这就是三维体积元 $\mathrm{d}V$，即超曲面元在超平面 $x^0=\mathrm{const}$ 上的投影.

4. 沿一个四维体积的积分；积分元是标量
\begin{equation}
  \mathrm{d}\Omega=\mathrm{d}x^0\mathrm{d}x^1\mathrm{d}x^2\mathrm{d}x^3=c\,\mathrm{d}t\,\mathrm{d}V.
\end{equation}
这个积分元是一标量：四维空间一部分的体积在坐标系转动时显然是不变的.\footnote{在积分变量 $x^0,x^1,x^2,x^3$ 变换到变量 $x'^0,x'^1,x'^2,x'^3$ 时，积分元 $\mathrm{d}\Omega$ 变为 $J\mathrm{d}\Omega'$，这里
\begin{equation*}
  \mathrm{d}\Omega'=\mathrm{d}x'^0\mathrm{d}x'^1\mathrm{d}x'^2\mathrm{d}x'^3
\end{equation*}
\begin{equation*}
  J=\frac{\partial(x'^0,x'^1,x'^2,x'^3)}{\partial(x^0,x^1,x^2,x^3)}
\end{equation*}
是该变换的雅可比行列式.对于形为 $x'^i=\alpha^i_k x^k$ 的线性变换，雅可比行列式 $J$ 就是行列式 $|\alpha^i_k|$ 并且对于坐标系的转动等于 1；这显示了 $\mathrm{d}\Omega$ 的不变性.}

类似三维矢量分析中的高斯定理和斯托克斯定理，有些定理使我们能做四维积分的变换.

沿一闭合超曲面的积分可以变换到沿包含在它里面的四维体积的积分，办法是用算符
\begin{equation}
  \mathrm{d}S_i\to\mathrm{d}\Omega\frac{\partial}{\partial x^i}
\end{equation}
代替积分元 $\mathrm{d}S_i$.例如，对于矢量 $A^i$ 的积分，我们有：
\begin{equation}
  \oint A^i\mathrm{d}S_i=\int\frac{\partial A^i}{\partial x^i}\mathrm{d}\Omega.
\end{equation}
这个公式是高斯定理的推广.

沿一个二维曲面的积分可以变换为“包含”它的超曲面的积分，办法是用算符
\begin{equation}
  \mathrm{d}f^{*}_{ik}\to\mathrm{d}S_i\frac{\partial}{\partial x^k}
  -\mathrm{d}S_k\frac{\partial}{\partial x^i}.
\end{equation}
代替积分元 $\mathrm{d}f^{*}_{ik}$.例如，对于反对称张量 $A^{ik}$ 的积分，我们有：
\begin{equation}
  \frac{1}{2}\oint A^{ik}\mathrm{d}f^{*}_{ik}
  =\frac{1}{2}\int\left(\mathrm{d}S_i\frac{\partial A^{ik}}{\partial x^k}
  -\mathrm{d}S_k\frac{\partial A^{ik}}{\partial x^i}\right)
  =\int\mathrm{d}S_i\frac{\partial A^{ik}}{\partial x^k}.
\end{equation}

沿一条四维闭合曲线的积分可以通过代换：
\begin{equation}
  \mathrm{d}x^i\to\mathrm{d}f^{ki}\frac{\partial\ \ }{\partial x^k}
\end{equation}
变换为“包含”它的曲面的积分.因此，对于一个矢量的积分，我们有：
\begin{equation}
  \oint A_i\mathrm{d}x^i
  =\int\mathrm{d}f^{ki}\frac{\partial A_i}{\partial x^k}
  =\frac{1}{2}\int\mathrm{d}f^{ki}
  \left(\frac{\partial A_k}{\partial x^i}-\frac{\partial A_i}{\partial x^k}\right),
\end{equation}
这是斯托克斯定理的推广.

\begin{xiti}

\noindent{\heiti 习题1}\quad 求一个对称四维张量 $A^{ik}$ 的分量在洛伦兹变换（6.1）下的变换法则.

解：将该张量的分量看做两个四维矢量分量的乘积，我们得到：
\begin{equation*}
  A^{00}=\frac{1}{1-\dfrac{V^2}{c^2}}
  \left(A^{\prime 00}+2\frac{V}{c}A^{\prime 01}
  +\frac{V^2}{c^2}A^{\prime 11}\right),
\end{equation*}
\begin{equation*}
  A^{11}=\frac{1}{1-\dfrac{V^2}{c^2}}
  \left(A^{\prime 11}+2\frac{V}{c}A^{\prime 01}
  +\frac{V^2}{c^2}A^{\prime 00}\right),
\end{equation*}
\begin{equation*}
  A^{22}=A^{\prime 22},\quad A^{23}=A^{\prime 23},\quad
  A^{12}=\frac{1}{\sqrt{1-\dfrac{V^2}{c^2}}}
  \left(A^{\prime 12}+\frac{V}{c}A^{\prime 02}\right),
\end{equation*}
\begin{equation*}
  A^{01}=\frac{1}{1-\dfrac{V^2}{c^2}}
  \left[A^{\prime 01}\left(1+\frac{V^2}{c^2}\right)
  +\frac{V}{c}A^{\prime 00}+\frac{V}{c}A^{\prime 11}\right],
\end{equation*}
\begin{equation*}
  A^{02}=\frac{1}{\sqrt{1-\dfrac{V^2}{c^2}}}
  \left(A^{\prime 02}+\frac{V}{c}A^{\prime 12}\right),
\end{equation*}
以及对于 $A^{33},A^{13}$ 和 $A^{03}$ 的类似公式.

\noindent{\heiti 习题2}\quad 对反对称张量 $A^{ik}$ 求解同样的问题.

解：因为坐标 $x^2$ 和 $x^3$ 不变，张量分量 $A^{23}$ 就不变，而分量 $A^{12},A^{13}$ 和 $A^{02},A^{03}$ 像 $x^1$ 和 $x^0$ 一样变换：
\begin{equation*}
  A^{23}=A^{\prime 23},\quad
  A^{12}=\frac{A^{\prime 12}+\dfrac{V}{c}A^{\prime 02}}
  {\sqrt{1-\dfrac{V^2}{c^2}}},\quad
  A^{02}=\frac{A^{\prime 02}+\dfrac{V}{c}A^{\prime 12}}
  {\sqrt{1-\dfrac{V^2}{c^2}}}.
\end{equation*}
类似地可得 $A^{13},A^{03}$.

对于 $x^0x^1$ 平面内二维坐标系的转动（它就是我们正在考虑的变换），分量 $A^{01}=-A^{10}$，$A^{00}=A^{11}=0$，构成一个2阶反对称张量，阶数2等于空间维数.因此（见式（6.9）后第2个脚注），这些分量在该变换下不变：
\begin{equation*}
  A^{01}=A^{\prime 01}.
\end{equation*}

\end{xiti}

\section{四维速度}

由普通的三维速度矢量，我们可以构造一个四维矢量.一个粒子的四维速度（四速度）是矢量
\begin{equation}
  u^i=\frac{\mathrm{d}x^i}{\mathrm{d}s}.
\end{equation}

为了求出它的分量，我们应注意，根据（3.1），
\begin{equation*}
  \mathrm{d}s=c\,\mathrm{d}t\sqrt{1-\frac{v^2}{c^2}},
\end{equation*}
其中 $v$ 为粒子的普通三维速度.因此，
\begin{equation*}
  u^1=\frac{\mathrm{d}x^1}{\mathrm{d}s}
  =\frac{\mathrm{d}x}{c\,\mathrm{d}t\sqrt{1-\dfrac{v^2}{c^2}}}
  =\frac{v_x}{c\sqrt{1-\dfrac{v^2}{c^2}}},
\end{equation*}
我们用同样的方法来求 $u^0,u^2,u^3$，结果我们得到：
\begin{equation}
  u^i=\left(\frac{1}{\sqrt{1-\dfrac{v^2}{c^2}}},
  \frac{v}{c\sqrt{1-\dfrac{v^2}{c^2}}}\right).
\end{equation}

应该注意，四维速度是一个无量纲量.

四维速度的分量并不彼此独立.注意到，$\mathrm{d}x_i\mathrm{d}x^i=\mathrm{d}s^2$，我们有
\begin{equation}
  u^iu_i=1.
\end{equation}

所以，就几何意义言之，我们可以说，$u^i$ 是与粒子世界线相切的一个四维单位矢量.

与四维速度的定义类似，二阶导数
\begin{equation*}
  w^i=\frac{\mathrm{d}^2x^i}{\mathrm{d}s^2}=\frac{\mathrm{d}u^i}{\mathrm{d}s}
\end{equation*}
可以称为四维加速度.微分（7.3），我们求得
\begin{equation}
  u_iw^i=0,
\end{equation}
即四维速度矢量和四维加速度矢量是相互正交的.

\begin{xiti}

\noindent{\heiti 习\ 题}\quad 确定相对论的匀加速运动，即在固有参考系中（每个时刻）加速度 $w$ 保持不变的直线运动.

解：在粒子速度 $v=0$ 的参考系中，四维加速度的分量 $w^i=(0,w/c^2,0,0)$（这里 $w$ 是普通三维加速度，指向沿 $x$ 轴）.相对论不变的匀加速条件必须表达为四维标量（在固有参考系中与 $w^2$ 一致）的常数性：
\begin{equation*}
  w^iw_i=\mathrm{const}\equiv-\frac{w^2}{c^4}.
\end{equation*}
在与之参照来观察运动的这个“固定”系中，写出 $w^iw_i$ 的表达式，得到方程
\begin{equation*}
  \frac{\mathrm{d}}{\mathrm{d}t}\frac{v}{\sqrt{1-\dfrac{v^2}{c^2}}}=w,\quad
  \text{或}\quad
  \frac{v}{\sqrt{1-\dfrac{v^2}{c^2}}}=wt+\mathrm{const}.
\end{equation*}
对于 $t=0$ 令 $v=0$，我们得到 $\mathrm{const}=0$，所以
\begin{equation*}
  v=\frac{wt}{\sqrt{1+\dfrac{w^2t^2}{c^2}}}.
\end{equation*}
再积分一次并对 $t=0$ 令 $x=0$，我们得到：
\begin{equation*}
  x=\frac{c^2}{w}\left(\sqrt{1+\frac{w^2t^2}{c^2}}-1\right).
\end{equation*}
对于 $wt\ll c$，这些公式过渡到经典表达式 $v=wt$，$x=wt^2/2$.对于 $wt\to\infty$，速度趋于恒定值 $c$.

一个匀加速粒子的固有时由如下积分给出
\begin{equation*}
  \int_0^t\sqrt{1-\frac{u^2}{c^2}}\,\mathrm{d}t
  =\frac{c}{w}\operatorname{arsinh}\frac{wt}{c}.
\end{equation*}
当 $t\to\infty$ 时，按照规律 $\dfrac{c}{w}\ln\dfrac{2wt}{c}$，它比 $t$ 增加慢得多.

\end{xiti}
